\section{Application Experiments and Interpretation}
The experimental unit is a stated invocation or application session, with a recorded workload, software revision, and network configuration. Report submitted, admitted, selected, completed, rejected, timed-out, and cancelled work separately where the path exposes those outcomes. Success rate uses an explicit denominator and deadline. Latency distributions should identify which requests they include, so faster completed requests do not hide a larger failure fraction. Repeated trials or independent seeds support variability estimates; a single exploratory run is reported as such.

For UAV sessions, record command completion latency, video start delay, useful frame delivery, recording retrieval completion, and status age at the operator. For multi-view work, additionally record the number of usable views, processing completion, and the analysis metric against labeled observations. The framework experiment first holds the recognition algorithm fixed. A separate application study is needed before attributing higher recognition accuracy to additional views.

For DI, report model preparation, prefill, time to first token, subsequent output rate, total completion, peak device memory, and bytes transferred between Providers. Compare cold and warm conditions separately. Evaluate a feasible single-Provider case alongside a supported partitioned case. If the distributed deployment exists because the model cannot fit one device, report that capacity benefit directly rather than inventing a speedup over an infeasible execution. Profile communication, serialization, cryptography, and model computation to locate the limiting cost.

Authorization evaluation should exercise both directions of service permission, not only successful decryption. Negative cases include an unpermitted caller, an unpermitted Provider, a validly signed message with the wrong protected value, and reuse after request state is closed. Grant experiments measure how an enrolled entity learns and applies a new permission. Withdrawal experiments measure update propagation, status validation, parameter/DKEY refresh, and rejection after the update is applied. Disconnected nodes and old ciphertext are separate cases with the limitations stated in Section~\ref{sec:abe-rationale}.

The collaboration ablations retain request-ID isolation and signature validation. Removing assigned-producer checks tests a wrong-role object within the same request; removing edge checks tests an undeclared input from an otherwise authorized participant. Replacement tests inject old work after reassignment. This isolates the extension's contribution instead of claiming to solve cross-request confusion that the invocation foundation already handles. Correctness tests require the expected rejection and absence of prohibited handler execution, not merely an error message in a log.

Framework reuse can be assessed by documenting which invocation, authorization, object-transfer, and lifecycle components are shared by UAV and DI, and which adapters remain application-specific. Source size alone is not a sufficient maintainability metric. Record the changes needed to add a service type and the contracts that must be supplied. This supports the framework's engineering rationale while leaving protocol correctness and performance to their own evidence.

Deployment records should include dependency versions, model hashes, topology, link parameters, clocks, retry policies, cache configuration, and key-management settings without publishing secrets. Separate application failures from startup, configuration, and collection failures. A failed setup is not evidence of a protocol's inability to complete a request. These records allow the final dissertation to distinguish supported claims, useful negative results, and configurations that remain untested.
